% !TeX root = ../main.tex
% ============================================================
%  sections/04b_leader.tex
% ============================================================

%-----------------------------------------------------------------
\subsection{Bilevel Programme}
\label{sec:model:bilevel}

The \CLSC design problem with the carbon tax as a leader decision
variable is modelled as
the following bilevel programme. The leader announces a carbon-tax index
$\taxhat$; the follower observes $\taxhat$ and selects the
cost-minimising supply chain design and operations. Denoting the
follower's optimal-response correspondence by $\Sol{\taxhat}$ and its
feasible region by $\setX$ (defined in \Cref{sec:model:follower}), the
complete programme is
\begin{subequations}\label{eq:leader}
\begin{alignat}{2}
\textup{(BL)}\quad
\lexmin_{\taxhat,\,\bz}\quad
  & \bigl(\Emis(\bz),\ \taxhat\bigr) \label{eq:leader_obj}\\
\text{s.t.}\quad
  & \taxhat\in\mathcal{T}^{\tau}, \label{eq:leader_tax}\\
  & \bz\in\Sol{\taxhat}, \label{eq:leader_react}
\end{alignat}
\end{subequations}
where $\Sol{\taxhat}$ is defined in~\eqref{eq:reaction_set} once the
follower's objective has been introduced in \Cref{sec:model:follower_obj}.
\Cref{sec:model:leader,sec:model:follower} present the two levels in
detail.

Bilevel programmes with integer lower-level variables are NP-hard in
general \citep{moore1990, jeroslow1985}; problem~\eqref{eq:leader}
inherits this hardness, since the follower's problem~\eqref{eq:follower}
is itself a \MILP for every fixed $\taxhat$.

%-----------------------------------------------------------------
\subsection{Regulator's Problem}
\label{sec:model:leader}

%- - - - - - - - - - - - - - - - - - - - - - - - - - - - - - - -
\subsubsection{Objective function}
\label{sec:model:leader_obj}

The regulator is a \emph{lexicographic minimiser}: it first minimises
the total emissions $\Emis(\bz)$ induced by the follower's response and,
among all tax indices that attain that minimum, selects the smallest one
(see~\eqref{eq:leader_obj}).

\begin{remark}[Lexicographic structure]\label{rem:lexicographic}
The lexicographic objective~\eqref{eq:leader_obj} minimises emissions
unconditionally, with the tax index acting solely as a tiebreaker. This
secondary criterion encodes a \emph{least-intervention principle}: the
regulator does not levy a tax heavier than the emissions target
requires. \Cref{sec:exp:tradeoff} shows this matters when the
emission--tax response saturates and the primary minimiser is attained
across a wide band of tax indices.
\end{remark}

%- - - - - - - - - - - - - - - - - - - - - - - - - - - - - - - -
\subsubsection{Feasibility constraints}
\label{sec:model:leader_constr}

Two constraints define the leader's feasible set. Constraint~\eqref{eq:leader_tax}
restricts the tax to the admissible grid $\mathcal{T}^{\tau}$
(\Cref{sec:model:variables}). Constraint~\eqref{eq:leader_react}
imposes bilevel feasibility: the supply chain operator's response must
be optimal given the announced tax. Under the optimistic
tie-break of \Cref{as:opt}, the leader selects, among all optimal
follower responses, the one with the lowest emissions.
